\thispagestyle{otFront}
\phantomsection\pdfbookmark[0]{Sommaire}{bm-sommaire}
\otPageTitre{Sommaire}{Terminale générale · Spécialité physique-chimie}
\vspace{-6mm}\sommaire\clearpage
\thispagestyle{otFront}
\otPageTitre{Parcours et ressources}{Compléments, entraînements et outils de révision}
\fontsize{11}{16}\selectfont
\otIntertitre{Approfondir}
\otLigneRessource{defis-debut}{\otNbDefis\ défis d'approfondissement}
\otLigneRessource{indices-defis}{Indications progressives pour les défis}
\otIntertitre{Préparer l'épreuve}
\otLigneRessource{epreuve-guide}{Format de l'épreuve, conseils et grille de rédaction}
\otLigneRessource{ece-guide}{Préparer l'évaluation des compétences expérimentales}
\otLigneRessource{fiches-experimentales}{Fiches expérimentales : verrerie, montages, mesures}
\otLigneRessource{sujet-1}{\otNbSujetsLettresMaj\ sujets type bac de 3 h 30}
\otIntertitre{Comprendre et corriger}
\otLigneRessource{corriges-debut}{Corrigés des \otNbExosChapitres\ exercices de chapitre}
\otLigneRessource{defi-cor-1-1}{Corrigés détaillés des \otNbDefis\ défis}
\otLigneRessource{sujet-1-cor}{Corrigés et barèmes des sujets type bac}
\otIntertitre{Se repérer}
\otLigneRessource{donnees-utiles}{Données utiles : constantes, tableau périodique, électronégativités}
\otLigneRessource{couverture-programme}{Repères du programme : notions et capacités exigibles}
\otLigneRessource{references}{Références officielles}
\clearpage
